%%%%%
%
\section{Extraction de "colonnes" cons\'ecutives d'un vecteur 1D de type \texttt{float}}

\itemstyle{Op\'erateur~:} \operfmt{Extr\_tabflt2flt}
\index{Extr\_tabflt2flt}

%%%%%OBJET%%%%
\itemstyle{Objet~:}
 Extrait une ou plusieurs �colonnes� cons\'ecutives d'un vecteur 1D
de type \texttt{float}, et les range dans un nouveau vecteur 1D
du m�me type.

%%%%%ARGUMENTS%%%%
\bigskip
%%%%%ENTREES%%%%
\begin{tabarg}
\multicolumn{3}{l}{\titretab {Entr\'ees}}\\
\hline
$X$ & Vecteur, \texttt{float[$d_x$]} & Param\`etre 1D d'entr\'ee (unit� indiff�rente)\\
$d_x$ & Coeff., \texttt{Int} & Taille  de la dimension (hors-temps) de $X$ (nombre de �colonnes�)\\
$d_z$ & Coeff., \texttt{Int} & Taille de la dimension (hors temps) du param�tre extrait (nombre de �colonnes�)\\
$C_{deb}$ & Coeff., \texttt{Int} & Indice de d�but de l'extraction (convention "C", d�but � z�ro)\\
\end{tabarg}

\medskip
%%%%%SORTIES%%%%
\begin{tabarg}
\multicolumn{3}{l}{\titretab {Sorties}}\\
\hline
$Z$ & Vecteur, \texttt{float[$d_z$]} & Sous-tableau (m\^eme unit\'e que $X$)\\
\end{tabarg}


%%%%%FORMULE%%%%
\itemstyle{Formule}

%%Extr_tabflt2flt
\[
Z_k = X_{k + C_{deb}}  \qquad k=0, d_z - 1
\]



%%%% COMMENTAIRE %%%%
\itemstyle{Commentaire~:}

Lorsque l'argument $d_{z}$ (Le nombre de �colonnes� � extraire) 
est �gal � 1, le vecteur r�sultat est un vecteur �1D� dont
la taille de la dimension est 1, ce qui est parfaitement
�quivalent � un vecteur �scalaire�.

%%%%%AUTEUR%%%%
\itemstyle{Auteur~:}
CNRM/GMEI/TRAMM

%%%%%REFERENCE%%%%
\itemstyle{R\'ef\'erences~:}
N�ant.
